\section{Introduction}
\label{sec:intro}

In~the~last twenty 
years 
a~plethora of new hadron states
have been discovered 
in decays of beauty hadrons 
to
%%
%% {\color{red}{\sout{charmonia,}}}
%% {\color{red}{
    charmonium,
%% }}
%%
including 
the~enigmatic $\chicone(3872)$~state~\cite{Choi:2003ue}, 
numerous pentaquark 
states in $\jpsi\proton$~\cite{LHCb-PAPER-2015-029,
LHCb-PAPER-2016-009,
LHCb-PAPER-2019-014,
LHCb-PAPER-2016-015,
LHCb-PAPER-2021-018}
and $\jpsi\Lz$~\cite{LHCb-PAPER-2020-039}~systems
as well as tetraquarks  
in 
the~$\psitwos\pip$~\cite{Choi:2007wga,
Mizuk:2009da,
Chilikin:2013tch,
LHCb-PAPER-2014-014,
LHCb-PAPER-2015-038},
$\jpsi\Pphi$~\cite{LHCb-PAPER-2016-018,
LHCb-PAPER-2016-019,
LHCb-PAPER-2020-035,
LHCb-PAPER-2020-044,
LHCb-PAPER-2022-040},
$\Peta_{\cquark}(1\PS)\pim$~\cite{LHCb-PAPER-2018-034},
$\jpsi\pip$~\cite{LHCb-PAPER-2018-043},
$\jpsi\Kp$~\cite{LHCb-PAPER-2020-044}
and 
$\jpsi\KS$~\cite{LHCb-PAPER-2022-040}~systems. 
Transitions among
%%
%%{\color{red}{\sout{charmonia or charmonia\nobreakdash-like}}}
%%{\color{red}{
    charmonium or charmonium\nobreakdash-like
%%}}
%%
states have been studied in 
beauty-hadron decays,  
including transitions with emission of 
one~photon~\cite{LHCb-PAPER-2013-024,
LHCb-PAPER-2014-008,
LHCb-PAPER-2017-011,
LHCb-PAPER-2021-003}, 
two pions~\cite{LHCb-PAPER-2013-001,
LHCb-PAPER-2015-015,
LHCb-PAPER-2020-009, 
LHCb-PAPER-2021-045},
$\Pphi$~\cite{LHCb-PAPER-2016-018,
LHCb-PAPER-2016-019,
LHCb-PAPER-2020-035,
LHCb-PAPER-2020-044,
LHCb-PAPER-2022-040},
$\Pomega$~\cite{BaBar:2010wfc,
LHCb-PAPER-2021-045} 
and $\Peta$~\cite{LHCb-PAPER-2021-047}~mesons.
The~transitions with emission of 
an~\Peta~meson 
have 
also 
been
studied in 
\mbox{$\decay{\epem}
{\jpsi\Peta}$}~processes~\cite{Wang:2012bgc,
BESIII:2020bgb}. 
In~general, studies of various hadronic transitions
in
%%
%% {\color{red}{\sout{the~charmonia and charmonia\nobreakdash-like}}}
%% {\color{red}{
the~charmonium and charmonium\nobreakdash-like
%% }}
%%
sectors 
can shed light
%%
%% {\color{red}{\sout{into}}}
%% {\color{red}{
    onto
%%}}
%%
the~internal 
structure of these particles, 
which is largely unknown for newly 
discovered hadronic states~\cite{Brambilla:2019esw,
Ali:2019roi,
Maiani:2022psl,
Brambilla:2022ura}.
%%

Transitions with an~emission of 
an~$\Peta^{\prime}$~meson in
%%
%% {\color{red}{\sout{the~charmonia and charmonia\nobreakdash-like}}}
%% {\color{red}{
    the~charmonium and charmonium\nobreakdash-like
%%}}
%%
systems have not yet 
been observed~\cite{Brambilla:2010cs,
Brambilla:2019esw,
PDG2022}. 
Since the~\etapr~meson may have 
a~glueball contribution~\cite{Rosner:1982ey,
Bramon:1997va,
Bramon:1997mf,
Novikov:1979uy,
Novikov:1979ux,
Novikov:1979va,
Kataev:1981aw,
Kataev:1981gr,
Thomas:2007uy,
Escribano:2007cd,
Escribano:2008rq,
Ambrosino:2009sc,
Yelton:2009aa,
PhysRevD00,
DiDonato:2011kr,
LHCb-PAPER-2014-056,
Andreichikov:2019ehn,
Andreichikov:2019uoa}, 
processes involving this particle are 
of particular interest~\cite{Close:1997wp,Close:2005iz}.
The~\betapk~decay\footnote{Inclusion of charge-conjugate 
states is implied throughout the~paper, 
unless otherwise 
stated.}
is a~good candidate 
to explore the~$\jpsi\etapr$~system in detail,
offering 
the~opportunity
to search 
for possible intermediate resonances. 
The~decay itself 
has never been observed
and 
an~upper limit on its branching fraction of 
\begin{equation*}
\BR(\betapk) < 8.8\times10^{-5}\ (90\,\%~\mathrm{CL})\,,
\end{equation*}
 was set 
by the~\belle collaboration~\cite{Belle:2006aam}.



This paper reports the~observation of 
the~\betapk~decay 
using 
  proton\nobreakdash-proton\,(\proton\proton) collision data
  collected by the~\lhcb experiment  
  at centre-of-mass energies of 
  7, 8, and 13\tev,
  corresponding 
  to a~total integrated luminosity of $9\invfb$.
%% 
The~measurement of its branching fraction  
normalised to 
the~well\nobreakdash-known 
branching fraction of the~$\bpsik$~decay~\cite{PDG2022}, 
\begin{equation}
\label{eq:ratio}
\mathcal{R}
\equiv
\dfrac{\BF(\betapk)}{\BF(\bpsik)} \,,
\end{equation}
%% and 
is performed using
the~\etaprhog~decay.
The~observation of the~signal 
is confirmed using 
the~\etapetapipi~decay mode, 
which is also used 
as a~cross\nobreakdash-check.


\section{Detector and simulation}
\label{sec:detector}


The \lhcb detector~\cite{LHCb-DP-2008-001,LHCb-DP-2014-002} 
is a~single\nobreakdash-arm forward
spectrometer covering 
the~\mbox{pseudorapidity} range 
\mbox{$2<\eta <5$},
designed for the study of particles containing \bquark or \cquark
quarks. The detector includes 
a~high\nobreakdash-precision tracking system
consisting of 
a~silicon\nobreakdash-strip vertex detector 
surrounding 
the~$\proton\proton$
interaction region~\cite{LHCb-DP-2014-001}, 
a~large\nobreakdash-area 
silicon\nobreakdash-strip detector located
upstream of a~dipole magnet with 
a~bending power of about~$4{\mathrm{\,Tm}}$, 
and three stations of 
silicon\nobreakdash-strip detectors 
and straw drift tubes~\cite{LHCb-DP-2013-003,
LHCb-DP-2017-001}
placed downstream of the~magnet.
The~tracking system provides a~measurement of 
the~momentum, \ptot, of charged particles with
a~relative uncertainty that varies 
from~0.5\% at low momentum 
to~1.0\% at 200\gevc.
The~minimum distance of a~track 
to a~primary $\proton\proton$~collision 
vertex\,(PV), the~impact parameter  %% \,(IP), 
is measured with a~resolution of~$(15+29/\pt)\mum$,
where \pt is the~component of 
the~momentum transverse to the~beam, in\,\gevc.
Different types of charged hadrons 
are distinguished using information
from two ring\nobreakdash-imaging 
Cherenkov detectors~\cite{LHCb-DP-2012-003}. 
Photons, electrons and hadrons are 
identified by a~calorimeter system consisting of
scintillating-pad and preshower detectors, 
an electromagnetic
and a~hadronic calorimeter. 
Muons are identified by 
a~system composed of alternating layers 
of iron and multiwire
proportional chambers~\cite{LHCb-DP-2012-002}.





The~online event selection 
is performed by
a~trigger~\cite{LHCb-DP-2012-004,
LHCb-DP-2019-001}, 
which consists of a~hardware stage, based on information 
from the~calorimeter and muon systems, 
followed by a~software stage, 
which applies a~full event reconstruction.
At~the~hardware trigger stage, events are required to have 
a~muon track with high transverse momentum or dimuon candidates 
in which the~product of 
the~\pt of the~muons 
has a~high value. 
In~the~software trigger, two oppositely charged muons are required 
to form a~good\nobreakdash-quality vertex that is significantly 
displaced from every~PV, with a~dimuon mass exceeding $2.7\gevcc$.

Simulated events are used to describe signal shapes 
and to compute the~efficiencies needed to determine 
the~branching fraction ratio.
In~the~simulation, $\proton\proton$~collisions are 
generated using \pythia~\cite{Sjostrand:2007gs} with 
a~specific \lhcb configuration~\cite{LHCb-PROC-2010-056}.
Decays of unstable particles are 
described by \evtgen~\cite{Lange:2001uf}, 
in which final\nobreakdash-state radiation is generated 
using \photos~\cite{davidson2015photos}.
The~interaction of the~generated particles with the~detector, 
and its response, are implemented using 
the~\geant toolkit~\cite{Allison:2006ve, *Agostinelli:2002hh} 
as described in Ref.~\cite{LHCb-PROC-2011-006}.
The~\pt~and rapidity\,($y$) 
spectra of the~$\Bu$~mesons in simulation 
are corrected
to match 
distributions in data.
The~correction factors are calculated 
by comparing the~observed
$\pt$ and $y$~spectra
for
a~high\nobreakdash-purity 
data sample of 
reconstructed \decay{\Bu}{\jpsi\Kp}~decays
with the~corresponding simulated samples. 
In~the~simulation, 
the~\mbox{$\decay{\Bu}{\jpsi\etapr\Kp}$}~decays 
are generated as phase\nobreakdash-space decays and corrected 
using  a~gradient boosted decision 
tree reweighting 
algorithm~\cite{Rogozhnikov:2016bdp}
to reproduce  the~$\jpsi\etapr$ and $\etapr\Kp$~mass
spectra 
observed in data.   
%% 
To~describe accurately the~variables used 
for kaon identification,  
the~corresponding  quantities in simulation 
are resampled according to values obtained from
calibration data samples of 
\mbox{$\decay{\Dstarp}
{\left( \decay{\Dz}{\Km\pip}\right)\pip}$}~decays~\cite{LHCb-DP-2018-001}. 
The~procedure accounts for correlations between 
the~variables associated
%%
%% {\color{red}{\sout{to}}}
%% {\color{red}{
    with
%% }}
%%
a~particular track, 
as well as the~dependence of 
the~kaon identification response
on
the~track's $\pt$ and  $\eta$~and  the~multiplicity 
of tracks in the~event. 
%%
To~account for imperfections in the~simulation of 
charged\nobreakdash-particle reconstruction, 
the~track reconstruction efficiency
is corrected using 
a~sample of \mbox{$\decay{\jpsi}{\mumu}$}~decays 
in data~\cite{LHCb-DP-2013-002}.
%%
Samples of 
the~\mbox{$\decay{\Bu}{\jpsi\Kstarp}$}~decays 
with 
\mbox{$\decay{\Kstarp}{\Kp\left(\decay{\piz}{\g\g}\right)}$}
are used to correct 
the~photon reconstruction efficiency  
in simulation~\cite{LHCb-PAPER-2012-022,
LHCb-PAPER-2012-053,
Govorkova:2015vqa,
Govorkova:2124605}.

\section{Event selection}
\label{sec:sel}

The \betapk candidates are reconstructed 
with $\etapr$~decays to
%%
%% {\color{red}{\sout{both}}}
%% {\color{red}{
    either
%%}}
%%
%% through 
%% either of the~two \etapr decays with
%% the \mbox{$\decay{\etapr}{\left(\decay{\Prho^0}{\pip\pim}\right)\g}$}
%% or 
%% \mbox{$\decay{\etapr}{\left(\decay{\Peta}{\g\g}\right)\pip\pim}$}
\mbox{$\left(\decay{\Prho^0}{\pip\pim}\right)\g$}
or 
\mbox{$\left(\decay{\Peta}{\g\g}\right)\pip\pim$}
final states.
%% 
%% Although the~latter
%% has a~larger branching
%% fraction, 
%%
The~difficulty of reconstructing photons 
in the~\mbox{$\decay{\Peta}{\g\g}$}~decay 
leads to a~sample with fewer events.
%% 
The~\bpsik~normalisation decay
is reconstructed using 
the~\mbox{$\decay{\psitwos}{\jpsi\pip\pim}$}~decay. 
In~both signal and normalisation  
channels, the~\jpsi~meson 
is reconstructed in its decay to 
two muons.
%% 
As~explained in detail below, 
an~initial 
loose selection is 
applied for both 
signal and normalisation channels. 
Subsequently, for the~\betapk~candidates,
where the~background level is large, 
a~multivariate estimator
%% is applied.
is used to select higher purity subset of candidates.
The~normalisation channel 
has a~high purity 
after 
the~initial selection, 
therefore no further selection 
steps are applied.

To reduce systematic uncertainties,  
the~initial 
selection criteria for both signal and normalisation channels 
are kept the~same whenever possible. 
The~selection criteria are chosen to be similar to those used in 
previous LHCb studies~\mbox{\cite{LHCb-PAPER-2012-022,
LHCb-PAPER-2012-053,
LHCb-PAPER-2013-024,
LHCb-PAPER-2014-008,
LHCb-PAPER-2014-056,
LHCb-PAPER-2021-003,
LHCb-PAPER-2021-047}}. 
%%
The~muon, pion and kaon candidates are identified by combining 
information from the~Cherenkov~detectors,
calorimeters and muon detectors~\cite{LHCb-PROC-2011-008} 
associated
%%
%% {\color{red}{\sout{to}}}
%% {\color{red}{
    with
%% }}
%%
the~reconstructed tracks. 
To~reduce the~combinatorial background, 
only tracks that 
are inconsistent with originating from any 
reconstructed PV in the~event are considered. 
%% 
The~transverse momentum of the~muon candidates is required to
be greater than 500\mevc and their momenta 
must
exceed 6\gevc. Pairs of oppositely\nobreakdash-charged muons 
consistent with originating from a~common vertex 
are combined to form \mbox{$\decay{\jpsi}{\mumu}$}~candidates.
The~reconstructed mass 
of the~muon pair is required 
to be between $3.056$~and $3.136\gevcc$.

Tracks that are consistent with 
the~pion or kaon hypotheses are required to 
have transverse momentum greater than~200\mevc. 
Photons are reconstructed from clusters 
in the~electromagnetic 
calorimeter, 
with transverse energy above~350\mev.
The~clusters must not be 
associated with reconstructed 
tracks~\cite{Terrier:691743,
LHCb-DP-2020-001}.
Photon identification is based on the~combined information 
from electromagnetic and hadronic calorimeters, 
scintillation pad, 
preshower detectors and the~tracking system. 

For the~reconstruction  of 
the~\etapetapipi~candidates, 
two photons are first combined 
to form an~$\Peta$~candidate. 
The~diphoton mass 
is restricted to lie within $\pm60\mevcc$ 
around 
the~known mass of the~$\Peta$~meson~\cite{PDG2022}. 
Each~$\Peta$~candidate is then combined with two 
oppositely\nobreakdash-charged pions 
to form an~\etapr~candidate. 
The~mass of the~combination is required to lie within 
$\pm45\mevcc$ around the~known mass of
the~\etapr~meson~\cite{PDG2022}.
For~an~\etapr~candidate  
reconstructed in 
the~\etaprhog~decay mode, 
the~$\Prho^0$~candidate is formed from
two 
oppositely\nobreakdash-charged 
pions. 
The~mass of this candidate is restricted to 
lie between 500 and 900\mevcc.
%% 
This asymmetric region around the~known mass of 
the~$\Prho^0$~meson~\cite{PDG2022}
takes into account the~shift 
of the~$\Prho^0$~line shape, 
due to the~electric\nobreakdash-dipole nature of
the~\mbox{$\decay{\etapr}
{\Prho^0\g}$}~transition~\cite{TASSO:1984xri,
Kolanoski:1984hu,
ARGUS:1990dcm,
CLEO:1991kmc,
CLEO:1992dty,
BESIII:2017kyd}.
A~photon 
is combined with 
the~$\Prho^0$~candidate
in order to form 
the~\etapr~candidate, 
whose mass 
is required to lie within 
$\pm30\mevcc$ of the~known \etapr~mass~\cite{PDG2022}.

Each selected \jpsi~candidate is combined with 
a~kaon track 
and either an~\etapr~candidate or 
two oppositely\nobreakdash-charged pions 
to form a~\Bp~candidate decaying 
into the~signal or normalisation modes, 
respectively. 
For the~\mbox{$\decay{\Bu}{\psitwos\Kp}$}~candidates
the~$\jpsi\pip\pim$~mass is required to be 
between $3.66$ and $3.71\gevcc$.
To~improve the~\Bp~meson mass resolution 
a~kinematic fit~\cite{Hulsbergen:2005pu} 
is performed,
which 
constrains 
the~masses of the~\jpsi, \etapr and \Peta~candidates 
to their known values~\cite{PDG2022}, 
and the~\Bp~candidates to originate from 
its associated PV.
The~decay time of 
the~\Bp~candidates is required to be greater 
than $100\mum/c$ to suppress 
the~large combinatorial 
background from tracks created in a~PV.

Further selection of
the~\mbox{$\decay{\Bu}{\jpsi\etapr\Kp}$}~decays 
is based 
on a~multivariate estimator,
in the~following referred to as 
the~multi\nobreakdash-layer 
perceptron\,({\sc{MLP}}) classifier.
The~classifier 
is based on an~artificial neural 
network algorithm~\cite{McCulloch,rosenblatt58}, 
configured with a~cross\nobreakdash-entropy cost 
estimator~\cite{Zhong:2011xm}.
It~reduces the~combinatorial background to a~low level 
while retaining a~high signal efficiency. 
Two~{\sc{MLP}}~classifiers are trained separately for 
the~two 
different \etapr~meson decay modes.
%%
%% {\color{red}{\sout{In~both cases}}}
%% {\color{blue}{  %% \sout{
%% {\color{red}{For both decay modes}}
%% the~same variables are included 
%% in the~classifiers
%% and 
%% these are related to the~reconstruction quality, 
%% kinematics and decay time of the~\Bp~candidates, 
%% kinematics of the~final\nobreakdash-state particles 
%% and a~variable that characterises the~kaon identification.
%% }}
%%
%% {\color{magenta}{
    The~list of variables
    used for classifiers includes
    the~\chisq of the~kinematic fit;
    transverse momenta of the~\etapr,
    kaon and pion candidates;
    pseudorapidities of pion and kaon candidates;
    transverse momentum of the~photon
    from the~\mbox{$\decay{\etapr}{\rhoz\g}$}~decay
    or miminal transverse momentum of
    photons from
    the~\mbox{$\decay{\etapr}{\left(\decay{\Peta}{\g\g}\right)\pip\pim}$}~decay;
    decay time  of the~\Bu~candidate; 
    variable related to the~quality of kaon
    identification~\cite{LHCb-PROC-2011-008,LHCb-DP-2012-003}
    and,
    for the~\mbox{$\decay{\etapr}{\rhoz\g}$}~decay,
    %%{\color{red}{\sout{the~cosine of the~\rhoz~polarization angle in the~\etapr~rest frame.}}}
    %% {\color{red}{
        cosine of the~angle between momenta of the~\pip
        and \etapr~candidates in the~rest frame of the~\rhoz~candidate.
   %% }}
%% }}
%%
The~classifiers are trained using simulated samples 
of \betapk~decays as signal proxy, while  
the~\betapk~candidates from data with mass
above 5.35\gevcc are used to represent the background. 
The~\mbox{$\decay{\Bu}{\jpsi\left(\etaprhog\right)\Kp}$}
candidates  
with the~\jpsi\pip\pim\Kp mass
%%
%% {\color{red}{\sout{
%%      within $\pm21\mevcc$ 
%%      around}}}
%% {\color{red}{
    consistent with
%% }}
%%
the~known mass of the~\Bp~meson 
are vetoed to avoid 
contamination from 
the~\mbox{$\decay{\Bp}{\jpsi\pip\pim\Kp}$}~decays 
with a~random photon added.

The~requirement on each of the~{\sc{MLP}} classifiers 
is chosen to maximise 
the~figure\nobreakdash-of\nobreakdash-merit 
defined as 
%% $\tfrac{S}{\sqrt{B+S}}$,
$S/\sqrt{B+S}$,
where $S$~represents 
the~expected signal yield,
and $B$~is the~expected background yield
within a~$\pm15\mevcc$ 
mass window centred
around the~known mass of 
the~$\Bu$~meson
%%
%% {\color{red}{
    and corresponding to
    approximately
    three times the~mass resolution
    on both sides of the~peak.
%% }}
%%
The~background yield is calculated 
from fits to data,
%%
%% {\color{red}{
    as described in Sec.~\ref{sec:sig},
%% }}
%%
while 
the~expected signal yield is estimated as $S=\varepsilon S_0$,
where $\varepsilon$ is the~efficiency of the~requirement 
on the~response of the~{\sc{MLP}}~classifier determined 
from simulation, 
and $S_0$~is the~signal yield obtained 
from the~fit to 
the~data,
%%
%% when no 
with a~loose 
requirement
applied
on the~response of the~{\sc{MLP}}~classifier.\footnote{The~optimal requirement is found to be
  largely independent on the~choice
  of the~normalisation point.}
%%
The~mass distributions 
for the~selected \mbox{$\decay{\Bu}
{\jpsi\etapr\Kp}$}~candidates 
are shown in Fig.~\ref{fig:signals}, 
where clear signal
peaks 
corresponding to \Bp~decays are seen in data for 
both \etapr decay modes. 



\section{Signal yield determination}
\label{sec:sig}

The signal yields are determined using an~extended unbinned 
maximum\nobreakdash-likelihood fit
%%
%% {\color{red}{
    to the~\mbox{$\jpsi\etapr\Kp$}~mass distributions
%% }}
%%
with a~two\nobreakdash-component function.
The~signal component for both cases is modelled 
by a~modified Gaussian function
%%
%% {\color{red}{
    that combines a~Gaussian core
%% }}
%%
with power\nobreakdash-law tails 
on both sides of the~distribution~\cite{Skwarnicki:1986xj,
LHCb-PAPER-2011-013}. 
The~background component 
is parameterised by a~second\nobreakdash-order 
positive polynomial function~\cite{karlin1953geometry}
for
%%
%% {\color{red}{\sout{\Bp~mesons reconstructed through}}}
%%
the~\etaprhog~decay mode 
and by an~exponential function in 
case of 
the~\etapetapipi~decay mode.
%%
%% {\color{red}{\sout{The~tail parameters of the~signal components are taken 
%% from simulation. 
%% The~detector resolution
%% is taken from simulation and 
%% }}}
%% {\color{red}{
    The~parameters of the~detector
    resolution
    function are taken from simulation,
    and the~width of the~Gaussian function is
    %% }} 
further 
corrected by a~scale factor, $s_{\Bu}$,
%%
that accounts for a~small discrepancy between data and 
simulation~\cite{LHCb-PAPER-2020-009,
LHCb-PAPER-2020-035,
LHCb-PAPER-2021-034,
LHCb-PAPER-2021-047,
LHCb-PAPER-2022-025}. 
To~account for the~uncertainty in 
the~tail parameters and resolution, 
the~fit is performed simultaneously 
for data and simulated samples, 
sharing the~same tail parameters, 
and allowing the~correction factor $s_{\Bu}$~to vary.
%%
%% {\color{red}{\sout{The~fit results are overlaid}}}
%% {\color{red}{
    The~resulting fit functions are overlaid
    with data distributions
%% }}
%%
in Fig.~\ref{fig:signals}
and the~signal yields are found to be 
\begin{subequations}
\begin{eqnarray}
 \left. N_{\decay{\Bu}{\jpsi\etapr\Kp}}
 \right|_{\decay{\etapr}{\Prho^0\g}} 
     & = & 
    \left(1.11\phantom{0}\pm0.11\phantom{0}\right)\times 10^3\,, \label{eq:ns_rhog}
  \\
 \left. 
  N_{\decay{\Bu}{\jpsi\etapr\Kp}}
  \right|_{ \decay{\etapr}{\Peta\pip\pim} }
   & = &  
   %% \ 228\phantom{.} \pm 28 \,, \ 
 \left( 0.228 \pm 0.028 \right) \times 10^3 \,,
   \label{eq:ns_etapipi}
\end{eqnarray}
\end{subequations}
where the~uncertainties are statistical only. 
The~resolution correction factors $s_{\Bu}$~are found to be 
$1.08 \pm 0.12$ 
and $1.03 \pm 0.16$ 
for the~\etaprhog~and 
\etapetapipi~samples, respectively.  
In both cases, 
the~statistical significance of the 
\betapk~signal is 
calculated using Wilks' theorem~\cite{Wilks:1938dza}
and found 
to exceed  $17$ and $12$~standard deviations
for the~\etaprhog~and 
\etapetapipi~samples, respectively.  
However, as the~signal yield is much lower 
for the~\etapr~meson 
decays to the~$\Peta\pip\pim$~final state, 
all~subsequent studies are performed using only 
the~\etaprhog~decay mode.

\begin{figure}[t]
	\setlength{\unitlength}{1mm}
	\centering
	\begin{picture}(160,60)
	\put(  0, 0){\includegraphics*[width=80mm]{Fig1a.pdf}}
    \put( 80, 0){\includegraphics*[width=80mm]{Fig1b.pdf}}
	\put( 37,-2){${m_{\jpsi\etapr\Kp}}$}
	\put( 62,-2){$\left[\!\gevcc\right]$}
	\put( -2,31){\rotatebox[]{90}{Candidates\,/\,(1.5\mevcc)}}
	\put(117,-1){${m_{\jpsi\etapr\Kp}}$}
	\put(142,-1){$\left[\!\gevcc\right]$}
	\put( 80,32){\rotatebox[]{90}{Candidates\,/\,(3\mevcc)}}
	%%	
    \put( 55,37){$\decay{\etapr}{\Prho^0\g}$}
    \put(130,37){$\decay{\etapr}{\Peta\pip\pim}$}
%%
 \put(12,43){\scriptsize
   $\begin{array}{cl}
    %%\!\bigplus\mkern-5mu & \text{data} 
    \!\bigplus\mkern-17mu\bullet & \text{Data}
   \\
   \begin{tikzpicture}[x=1mm,y=1mm]\filldraw[fill=red!35!white,draw=red,thick]  (0,0) rectangle (5,3);\end{tikzpicture}  
   & \decay{\Bu}{\jpsi\etapr\Kp} 
   \\
   {\color[RGB]{85,83,246}{\hdashrule[0.0ex][x]{5mm}{1.0pt}{2.0mm 0.3mm}}}
   & \text{Background}
   \\ 
   {\color[RGB]{255,153,51} {\rule{5mm}{2.0pt}}}
   & \text{Total}
   \end{array}$ }
%% 
\put(92,43){\scriptsize
   $\begin{array}{cl}
    %% \!\bigplus\mkern-5mu & \text{data}  
    \!\bigplus\mkern-17mu\bullet & \text{Data}
   \\
   \begin{tikzpicture}[x=1mm,y=1mm]\filldraw[fill=red!35!white,draw=red,thick]  (0,0) rectangle (5,3);\end{tikzpicture}  
   & \decay{\Bu}{\jpsi\etapr\Kp} 
   \\
   {\color[RGB]{85,83,246}{\hdashrule[0.0ex][x]{5mm}{1.0pt}{2.0mm 0.3mm}}}
   & \text{Background}
   \\ 
   {\color[RGB]{255,153,51} {\rule{5mm}{2.0pt}}}
   & \text{Total}
   \end{array}$ } 
   %%  
   \put( 58,45){$\begin{array}{l}\lhcb\\ 9\invfb\end{array}$}
   \put(138,45){$\begin{array}{l}\lhcb\\ 9\invfb\end{array}$}
 %% 
	\end{picture}
	\caption { \small %captions should be a little bit smaller than main text
     Mass distributions 
     for selected \betapk~candidates 
     with $\etapr$~decays 
     to (left)~$\rhoz\g$
     and (right)~$\Peta\pip\pim$~final states. 
     %% are reconstructed using 
     %% (left)~\etaprhog~and 
     %% (right)~\etapetapipi~decay modes.     
    %% 
     %% {\color{red}{\sout{The results of the~fits,
     %%      described in the~text, 
     %%      are overlaid.}}}
     %% {\color{red}{
         The~resulting fit functions
         are overlaid with data distributions.
     %% }}
        }
	\label{fig:signals}
\end{figure}


The~background\nobreakdash-subtracted  
$\psietap$, 
$\etapk$,
and $\psik$~mass spectra 
from 
the~\mbox{$\decay{\Bu}
{\jpsi\left(\decay{\etapr}{\Prho^0\g}\right)\Kp}$}~decays
are shown in Figs.~\ref{fig:splots}(a-c),
where 
the~\sPlot~technique~\cite{Pivk:2004ty} 
based on the~fit results
is used for background subtraction. 
The~$\psietap$, 
$\etapk$,
and $\psik$~masses are calculated 
using a~kinematic fit 
with \jpsi, \etapr and \Bp mass constraints 
and a~PV constraint applied~\cite{Hulsbergen:2005pu}. 
While for the~$\jpsi\Kp$~mass 
the~distribution largely agrees with 
the~shape expected from the~phase\nobreakdash-space model, 
for the~low\nobreakdash-mass 
region of the~$\etapk$~mass spectrum 
and the~high\nobreakdash-mass 
region 
of the~$\psietap$~mass spectrum 
a~striking difference from the~phase\nobreakdash-space model 
is observed.
These differences 
are 
potentially 
due to 
contributions 
from decays via intermediate heavy 
excited strange mesons,
such as 
$\PK^{*}_{0}(1430)^{+}$, 
$\PK^{*}_{2}(1430)^{+}$ or 
$\PK^{*}(1680)^{+}$~mesons, 
decaying into 
the~$\etapr\Kp$~final state.
The~decays 
of the~\Bp~mesons into a~\jpsi~meson 
and heavy excited strange mesons 
have been
studied 
in 
Refs.~\cite{LHCb-PAPER-2016-018,
LHCb-PAPER-2016-019,
LHCb-PAPER-2020-044}.
%%
%% {\color{red}{\sout{Large~contribution from the~decays with 
%% intermediate excited kaons 
%% could also explain}}}
%% {\color{red}{
    The~decays via intermediate excited kaons
    also contribute to
    %% }}
the~higher mass region of the~$\jpsi\etapr$~mass spectrum,
as shown in Fig.~\ref{fig:splots}(d).
%%
\begin{figure}[t]
	\setlength{\unitlength}{1mm}
	\centering
	\begin{picture}(150,120)
	\put( 0,55){\includegraphics*[width=70mm]{Fig2_m23.pdf}}
	\put(72,55){\includegraphics*[width=70mm]{Fig2_m13.pdf}}
	\put( 0, 0){\includegraphics*[width=70mm]{Fig2_m12.pdf}}
	\put(72, 0){\includegraphics*[width=70mm]{Fig2_2d.pdf}}
    %% 
	\put( 35,52){${m_{\etapr\Kp}}$}
	\put(110,52){${m_{\jpsi\etapr}}$}
	\put( 35,-1){${m_{\jpsi\Kp}}$}
        \put(102,-1){${m_{\jpsi\etapr}}$}
	%%
	\put( 52, 52){$\left[\!\gevcc\right]$}
	\put(125, 52){$\left[\!\gevcc\right]$}
	\put( 52, -1){$\left[\!\gevcc\right]$}
        \put(120, -1){$\left[\gevcc\right]$}
        %%
	\put( -2, 86){\rotatebox[]{90}
        {$\frac{1}{N}\frac{\mathrm{d}N}
        {\mathrm{d}\,m_{\etapr\Kp}}\,\left[\frac{1}{\!\gevcc}\right]$}}
        %%
  	\put( 68, 86){\rotatebox[]{90}
        {$\frac{1}{N}\frac{\mathrm{d}N}
        {\mathrm{d}\,m_{\jpsi\etapr}}\,\left[\frac{1}
        {\!\gevcc}\right]$}}
        %%
        \put( -2, 30){\rotatebox[]{90}
        {$\frac{1}{N}\frac{\mathrm{d}N}
        {\mathrm{d}\,m_{\jpsi\Kp}}\,\left[\frac{1}
        {\!\gevcc}\right]$}}
        %%
        \put( 68, 30){\rotatebox[]{90}
        {$m_{\etapr\Kp}~~~\left[\!\gevcc\right]$}}
	%%  
       \put(141, 26){\rotatebox[]{90}
        {$\frac{1}{N}\frac{\mathrm{d}^2N}
        {\mathrm{d}\,m_{\etapr\Kp}\,
        \mathrm{d}\,m_{\jpsi\etapr}}
        \,\left[\frac{1}
        {\!\gev^2/c^4}\right]$}}
        
    \put(124,95){$\begin{array}{l}\lhcb\\ 9\invfb\end{array}$}
    \put( 51,95){$\begin{array}{l}\lhcb\\ 9\invfb\end{array}$}
    \put(118,40){$\begin{array}{l}\lhcb\\ 9\invfb\end{array}$}
    \put( 51,40){$\begin{array}{l}\lhcb\\ 9\invfb\end{array}$}
    %% 
    \put(12,  98){(a)}
    \put(85,  98){(b)}
    \put(12,  43){(c)}
    \put(85,  43){(d)}
    \put(22,90){\small
   $\begin{array}{cl}
    \!\bigplus\mkern-18mu{\scriptsize{\bullet}} & \text{Data}
   \\
   {\color[RGB]{0,0,255} {\rule{5mm}{1.0pt}}}
   & \text{Phase~space}
   \end{array}$ }
	\end{picture}
	\caption { \small 
        Normalised background\protect\nobreakdash-subtracted 
        (a)~$\etapk$,
        (b)~$\psietap$,
        (c)~$\psik$~mass spectra 
        and 
        (d)~two\protect\nobreakdash-dimensional mass distribution of
        $\etapr\Kp$ vs $\jpsi\etapr$
        from 
        the~\mbox{$\decay{\Bu}{\jpsi\etapr\Kp}$}~decays.
        %%
        Superimposed curves are the expectations 
        from a~phase\protect\nobreakdash-space model.
        }
	\label{fig:splots}
\end{figure}
%%
%%{\color{red}{\sout{While the~higher mass region  
%% of the~$\jpsi\etapr$~mass spectrum is 
%% presumably dominated by the~contributions
%% from the~decays via intermediate  
%% heavy excited kaons,}}} 
The~$\jpsi\etapr$~mass region below~$4.7\gevcc$
is explicitly inspected for 
possible contributions 
from decays via excited
%%
%% {\color{red}{\sout{charmonia or charmonium\nobreakdash-like}}}
%% {\color{red}{
    charmonium or charmonium\nobreakdash-like
%% }}
%%
states into the~$\jpsi\etapr$~final state. 
Fits~to the~background\nobreakdash-subtracted 
\psietap~mass distribution are performed 
in individual mass windows,
corresponding to 
the~well\nobreakdash-established 
%% resonances, 
%% namely 
$\Ppsi(4160)$, $\Ppsi(4230)$,
%% {\color{red}{\sout{$\Ppsi(3460)$,}}}
%% {\color{red}{
    $\Ppsi(4360)$,
%% }}
$\Ppsi(4415)$ 
and~$\Ppsi(4430)$~resonances~\cite{PDG2022}. 
For each fit, the~resonance shape is parameterised
with a~relativistic Breit\nobreakdash--Wigner function 
convoluted with a~mass 
resolution function.
%%
The~non\nobreakdash-resonant contribution is modelled 
by a~first order
%% {\color{red}{
    positive
%% }}
polynomial function.
The~known masses and widths of the resonances~\cite{PDG2022} 
are introduced in the~fits as Gaussian constraints 
on the~corresponding parameters.
%%
%% {\color{red}{\sout{The~resolution function is obtained using 
%% simulation as a~function of 
%% the~\psietap mass.}}}
%% {\color{red}{
    The~resolution function
    is modelled by the~modified Gaussian
    function with parameters 
    obtained using 
    simulation as a~function of 
    the~\psietap mass.
%% }}
%%
No~statistically significant signals 
are observed 
for the~\mbox{$\decay{\Bp}{\jpsi\etapr\Kp}$}~decays via 
intermediate 
resonances, listed above.
%% $\Ppsi(4160)$, $\Ppsi(4230)$,
%% $\Ppsi(3460)$, $\Ppsi(4415)$ or~$\Ppsi(4430)$~states
%% are observed. 
%%
%% {\color{red}{
    To~probe the contribution
    of the~resonances with higher masses,
    more advanced fit techniques
    accounnting for complicated background shape
    and  the~distortion of the~signal Breit\nobreakdash--Wigner
    shape are required.
%% }}



For the~determination of the~resonant structure
of the~\mbox{$\decay{\Bu}{\jpsi\etapr\Kp}$}~decay,
a~full amplitude analysis, 
similar to those used in 
Refs.~\cite{LHCb-PAPER-2016-019,
LHCb-PAPER-2020-044}, 
is required.
Large signal yields and 
low background levels 
are important prerequisites 
for such analysis. 
The~relatively large level of 
combinatorial background 
for 
the~\mbox{$\decay{\Bu}{\jpsi\etapr\Kp}$}~signal 
decays with the~\etapr~meson 
reconstructed via 
the~$\etaprhog$~decay mode, 
makes it difficult to carry out an~amplitude analysis.
With a~larger data sample, 
expected from 
future 
data\nobreakdash-taking periods,  
it will be possible to 
perform the~full amplitude 
analysis using 
the~\mbox{$\decay{\Bu}{\jpsi\etapr\Kp}$}~signal 
decays with the~\etapr~meson 
reconstructed via 
the~$\etapetapipi$~decay mode,
where the~combinatorial background 
is smaller. 




\section{Normalisation channel}
\label{sec:norm}

The \jpsi\pip\pim\Kp mass distribution 
for selected $\decay{\Bu}{\left(\decay{\psitwos}
{\jpsi\pip\pim}\right)\Kp}$~candidates
with $\jpsi\pip\pim$~mass between 3.66~and 3.71\gevcc 
is shown in Fig.~\ref{fig:norm_intermediates}(left). 
%%
An~extended unbinned 
maximum\nobreakdash-likelihood 
fit is performed,
where 
the~signal component
is modelled by 
%% {\color{red}{\sout{a}}}
the~modified Gaussian function
%%
%% {\color{red}{\sout{
%% with 
%% power\nobreakdash-law tails}}}
and the~background 
is described by 
a~first order
%% {\color{red}{
positive polynomial
%% l}}
function. 
The~result of this~fit is used to 
obtain the~background\nobreakdash-subtracted
$\jpsi\pip\pim$~mass distribution 
from \mbox{$\decay{\Bu}{\jpsi\pip\pim\Kp}$}~decays. 
This~distribution is shown in
Fig.~\ref{fig:norm_intermediates}(right). 
%%
The~yield of 
the~\mbox{$\decay{\Bu}
{\left(\decay{\psitwos}
{\jpsi\pip\pim}\right)\Kp}$}~signal
candidates is determined 
using an~unbinned fit to this distribution with 
a~two\nobreakdash-component function.
The~component corresponding 
to the~\mbox{$\decay{\Bu}{\psitwos\Kp}$}~decays
is parameterised 
by
%% {\color{red}{\sout{a}}}
the~modified Gaussian function.
%% with power\nobreakdash-law tails.
The~component describing
the~\mbox{$\decay{\Bu}{\jpsi\pip\pim\Kp}$}~decays
without intermediate \psitwos~state 
is modelled by
a~phase\nobreakdash-space 
function,\footnote{%%
The~phase\nobreakdash-space mass distribution of 
a~$k$\nobreakdash-body combination of
particles from an~$n$\nobreakdash-body decay 
is approximated by 
$\Phi_{k,n}(x) \propto 
x_{\ast}^{ (3k-5)/2}\left(1-x_{\ast}\right)^{3(n-k)/2-1}$,
where 
%%
$x_{\ast}\equiv 
( x-x_{\mathrm{min}} ) /(x_{\mathrm{max}}-x_{\mathrm{min}})$, 
and 
$x_{\mathrm{min}}$, $x_{\mathrm{max}}$ denote 
the~minimal   and maximal values of $x$, 
respectively~\cite{Byckling}.
Here, $k=3$ and $n=4$ are used.
} 
%% describing 
%% a~three\nobreakdash-body out of the~four\nobreakdash-body 
%% final state~\cite{Byckling},
modified by 
a~first order positive
polynomial function.
%%
The~tail and resolution parameters 
for the~signal component 
are taken from simulation
with the~resolution 
further corrected 
by a~scale factor, $s_{\psitwos}$, 
that accounts for a~small discrepancy 
between data and simulation~\cite{LHCb-PAPER-2020-008,
  LHCb-PAPER-2020-009}.
%%
%% {\color{red}{\sout{To~account for 
%% the~uncertainty in the~tail and resolution 
%% parameters,
%% the~fit is performed simultaneously
%% to
%% data and simulated samples, 
%% sharing the~same tail parameters, 
%% and allowing for the~correction 
%% factor $s_{\psitwos}$~to vary.}}}
%%
%% {\color{red}{
    The~fit is performed
      simultaneously to
      data and simulated samples,
      as for the~signal~mode
      described in Sec.~\ref{sec:sig}.
%% }}
%% 
From this fit the~number of 
\mbox{$\decay{\Bu}{\psitwos\Kp}$}~decays 
is found to be
\begin{equation}\label{eq:nn}
    %% N_{\bpsik} = 121398\pm141\,, 
    N_{\decay{\Bu}{\psitwos\Kp}} = \left( 121.40 \pm 0.14\right)\times 10^{3}\,, 
\end{equation}
where the uncertainty is statistical only. 
The~correction factor 
$s_{\psitwos}$~is found to be $1.057 \pm 0.008$, which is
in good~agreement with results from 
Refs.~\cite{LHCb-PAPER-2020-008,
LHCb-PAPER-2020-009,
LHCb-PAPER-2020-035,
LHCb-PAPER-2021-034,
LHCb-PAPER-2022-025,
LHCb-PAPER-2022-049}. 

\begin{figure}[t]
	\setlength{\unitlength}{1mm}
	\centering
	\begin{picture}(160,60)
	\put(0 , 0){\includegraphics*[width=80mm]{Fig3a.pdf}}
	\put(80, 0){\includegraphics*[width=80mm]{Fig3b.pdf}}
	\put( 35,-2){${m_{\psipipi\Kp}}$}
	\put(115,-2){${m_{\psipipi}}$}
	
        \put( 60,-2){$\left[\!\gevcc\right]$}
	\put(140,-2){$\left[\!\gevcc\right]$}
	\put( 2,29){\rotatebox[]{90}{Candidates\,/\,(1.5\mevcc)}}
	\put(82,34){\rotatebox[]{90}{Yield\,/\,(0.5\mevcc)}}
	%%  
    \put( 59,42){$\begin{array}{l}\lhcb\\ 9\invfb\end{array}$}
    \put(139,42){$\begin{array}{l}\lhcb\\ 9\invfb\end{array}$}
        %%
 \put(15,42){\scriptsize
   $\begin{array}{cl}
    \!\bigplus\mkern-17mu\bullet & \text{Data} 
   \\
   \begin{tikzpicture}[x=1mm,y=1mm]\filldraw[fill=red!35!white,draw=red,thick]  (0,0) rectangle (5,3);\end{tikzpicture}  
   & \decay{\Bu}{\jpsi\pip\pim\Kp} 
   \\
   {\color[RGB]{85,83,246}{\hdashrule[0.0ex][x]{5mm}{1.0pt}{2.0mm 0.3mm}}}
   & \text{Background}
   \\ 
   {\color[RGB]{255,153,51} {\rule{5mm}{2.0pt}}}
   & \text{Total}
   \end{array}$ }
%% 
\put(94,42){\scriptsize
   $\begin{array}{cl}
     \!\bigplus\mkern-17mu\bullet & \text{Data}
   \\
   \begin{tikzpicture}[x=1mm,y=1mm]\filldraw[fill=red!35!white,draw=red,thick]  (0,0) rectangle (5,3);\end{tikzpicture}  
   & 
   \decay{\Bu}{\psitwos\Kp}
   \\
   {\color[RGB]{85,83,246}{\hdashrule[0.0ex][x]{5mm}{1.0pt}{2.0mm 0.3mm}}}
   & \text{Non-res.} 
   \\ 
   {\color[RGB]{255,153,51} {\rule{5mm}{2.0pt}}}
   & \text{Total}
   \end{array}$ } 

        
	\end{picture}
	\caption { \small
    Left:~mass distributions 
    for selected \mbox{$\decay{\Bu}{\left(\decay{\psitwos}{\jpsi\pip\pim}\right)\Kp}$} decays, with 
    $\jpsi\pip\pim$~mass between 3.66~and~3.71\gevcc.
    Right:~background-subtracted  
    $\jpsi\pip\pim$~mass distribution 
    from selected \Bu~decays.  
    %% 
    %% {\color{red}{\sout{Results of the fits,
    %%       described in the~text, are overlaid.}}}
    %% {\color{red}{
    The~resulting fit functions
          are overlaid
          with data distributions.
    %% }}
    }
	\label{fig:norm_intermediates}
\end{figure}




\section{Branching 
fraction ratio computation}
\label{sec:bf}

The~ratio of the~branching fractions $\mathcal{R}$, 
defined in Eq.~\eqref{eq:ratio},
is calculated as 
\begin{equation*}
\mathcal{R}
= 
\dfrac{N_{\decay{\Bu}{\jpsi\etapr\Kp}}}
      {N_{\decay{\Bu}{\psitwos\Kp}}}
\times
\dfrac{\BF(\psitwos\to\jpsi\pip\pim)}{\BF(\etaprhog)}
\times 
%% \dfrac{\epsilon_{\bpsik}}{\epsilon_{\betapk}} 
\dfrac{\epsilon_{\decay{\Bu}{\psitwos\Kp}}}
      {\epsilon_{\decay{\Bu}{\jpsi\etapr\Kp}}} 
\,,
\end{equation*}
where 
$N_{\decay{\Bu}{\jpsi\etapr\Kp}}$
and 
$N_{\decay{\Bu}{\psitwos\Kp}}$
%%
are the~yields 
from Eqs.~\eqref{eq:ns_rhog} and~\eqref{eq:nn},  
and 
$\epsilon_{\decay{\Bu}{\jpsi\etapr\Kp}}$ 
and 
$\epsilon_{\decay{\Bu}{\psitwos\Kp}}$ 
are the~efficiencies
to reconstruct the~observed final states.
 The~efficiencies are the products of detector acceptance, 
reconstruction, selection 
 and trigger efficiencies, 
and 
 are calculated using simulated samples, 
 calibrated to match the data as 
 described in Sec.~\ref{sec:detector}.
 %%
 %% {\color{red}{
     The~branching fractions for
    the~\mbox{\decay{\psitwos}{\jpsi\pip\pim}}
    and \mbox{\decay{\etapr}{\rhoz\g}}~decays,
    \mbox{$\BR\left(\decay{\psitwos}{\jpsi\pip\pim}\right)= \left( 34.68 \pm 0.30\right)\%$}
    and
    \mbox{$\BR\left(\decay{\etapr}{\rhoz\g}\right)= \left( 29.4 \pm 0.4 \right)\%$},
    are taken from Ref.~\cite{PDG2022}.
%% }}
 %%
The~ratio of branching fractions $\mathcal{R}$ is found to be
\begin{equation*}\label{eq:R}
\dfrac{\BF(\betapk)}{\BF(\bpsik)} = 
\left(4.91\pm0.47\right)\times10^{-2},
\end{equation*}
where the uncertainty is statistical only. 
%%

As a~cross\nobreakdash-check, the~ratio 
$\mathcal{R}$~is also calculated for 
the~\mbox{$\decay{\etapr}{\Peta\pip\pim}$},
using the~signal yields 
from Eqs.~\eqref{eq:ns_etapipi} 
and~\eqref{eq:nn}.
The~ratio 
of branching fractions
$\mathcal{R}$ is found to be 
\mbox{$\left(4.98\pm0.61\right)\times10^{-2}$}, 
where the~uncertainty is statistical only. 
This value  is in 
good agreement with 
that 
calculated for the~\etaprhog~case.

\section{Systematic uncertainties}


%%{\color{red}{
The~signal and normalisation channels share 
the~same set of final\nobreakdash-state charged particles, 
and the~same trigger and preselection requirements
are applied to both.
%%
This allows 
many 
systematic uncertainties 
to cancel in the~ratio~$\mathcal{R}$. 
The~remaining nonnegligible 
contributions are listed in Table~\ref{tab:syst}.

%%
%% {\color{red}{\sout{A~large source 
%% of systematic uncertainties is}}}
%% {\color{red}{
Several systematic uncertainties are
%% }}
associated 
%%
%% {\color{red}{\sout{to}}}
%% {\color{red}{
    with
%% }}
%%
the~corrections applied to the~simulation. 
The~finite size of
the~$\decay{\Bp}{\jpsi\Kp}$ signal sample 
used for correction 
of the~simulated
transverse momentum and rapidity spectra 
of \Bp~mesons 
induces 
an~uncertainty
on the~\Bu~meson \pt and $y$~spectra.
In~turn, 
this uncertainty 
induces 
small changes in the~ratio of efficiencies. 
The~corresponding spread of these changes amounts to $0.1\%$ 
and is taken as the~systematic uncertainty
related to the~\Bu~meson kinematic.



The decay model corrections for 
the~\betapk~decay are obtained 
using 
%% the~gradient boosted 
%% tree reweighting algorithm~\cite{Rogozhnikov:2016bdp}.
the~algorithm described in Ref.~\cite{Rogozhnikov:2016bdp}.
The~systematic uncertainty related to the~correction 
method is estimated by varying 
%% the~hyperparameters  
%% of the~regression trees ensemble. 
the~configuration parameters of the~algorithm.
%%
The~largest deviation of the~efficiency value 
from the~baseline tuning is found to be $1.1\%$, 
which is assigned as systematic uncertainty 
associated with the~\Bu~decay model. 

There are residual differences in 
the~reconstruction efficiency 
of charged\nobreakdash-particle tracks that 
do not cancel completely in the~ratio 
of total efficiencies 
given 
the~slightly~different kinematic distributions 
of the~final\nobreakdash-state particles.
The~track\nobreakdash-finding efficiencies 
obtained from simulated samples are corrected using
calibration channels~\cite{LHCb-DP-2013-002}.
The~uncertainties related to~the~efficiency 
correction factors are propagated to the~ratios of 
the~total efficiencies using pseudoexperiments, 
and are found to be 0.7\%. 
This value is taken as the~systematic uncertainty 
due to the~tracking efficiency calibration.

Differences
%%
%% {\color{red}{\sout{on}}}
%% {\color{red}{
    in
%% }}
%%
the~photon reconstruction efficiencies 
between data and simulation 
are studied using a~large sample of 
\mbox{$\decay{\Bu}{\jpsi\Kstarp}$}~decays, 
reconstructed using 
the~\mbox{$\decay{\Kstarp}{\Kp\left(\decay{\piz}{\g\g}\right)}$}
decay mode~\cite{LHCb-PAPER-2012-022,
LHCb-PAPER-2012-053,
Govorkova:2015vqa,
Govorkova:2124605}.
%% 
%% {\color{red}{
The~uncertainty due to  
the~finite size of the~sample
is propagated to the~ratio of 
the~total efficiencies 
using pseudoexperiments
and is found
to be less 
than~1.0\%.
The~uncertainty due to the~accuracy 
of the~\mbox{$\decay{\Bu}{\jpsi\Kstarp}$}~branching 
fraction~\cite{PDG2022}
is 3.5\%.
%% }}
%% The~finite size of the~sample and 
%% the~accuracy 
%% of the~\mbox{$\decay{\Bu}{\jpsi\Kstarp}$}~branching 
%% fraction lead to corresponding systematic uncertainties. 
%% The~latter is independent of the~photon energy 
%% and it is known to be~3.5\%. 
%% The~former, depending on the~photon energy, 
%% is propagated to the~ratio of 
%% the~total efficiencies 
%% using pseudoexperiments
%% and is found
%% to be less 
%% than~1.0\%.
These two values are added in quadrature to obtain 
a~systematic uncertainty related to 
the~photon reconstruction of~3.6\%.

The~kaon identification variable used for the~{\sc{MLP}}~estimator 
is drawn from calibration data samples
and has a~dependence on the~particle
kinematics and track multiplicity. 
Systematic uncertainties in this procedure arise from 
the~limited
size 
of both the~simulation 
and calibration samples, 
and the~modelling of the~particle identification variable. 
The~limitations due to
the~size of the~simulation 
and calibration samples 
are evaluated by 
using bootstrapping
techniques~\cite{efron:1979,efron:1993}, 
creating  
multiple samples
and repeating the~procedure 
for
%% {\color{red}{\sout{each sample}}}
%% {\color{red}{
    each of these.
%% }}
The~impact of potential mismodelling of the~kaon identification  
variable is evaluated by describing 
the~corresponding distributions using 
density estimates with different kernel 
widths~\cite{LHCb-DP-2018-001,Poluektov:2014rxa}. 
For~each of these cases, alternative efficiency maps are 
produced to determine the~associated uncertainties. 
A~systematic uncertainty of 2.8\% is assigned from 
the~observed differences with alternative efficiency maps. 

A~systematic uncertainty  related 
to the~knowledge of the~trigger efficiencies 
has been previously studied using large 
samples of \mbox{$\decay{\Bu}
{\left(\decay{\jpsi}{\mumu}\right)\Kp}$} 
and \mbox{$\decay{\Bu}
{\left(\decay{\psitwos}{\mumu}\right)\Kp}$}~decays
by comparing the 
ratios of the trigger efficiencies in data 
and simulation~\cite{LHCb-PAPER-2012-010}.
Based on this comparison, a~relative uncertainty 
of 1.1\% is assigned.

The~remaining inconsistency between data and simulation, 
not covered by the~corrections discussed 
in Sec.~\ref{sec:detector}, 
is estimated by varying 
the~requirement on the~response
of the~{\sc{MLP}} classifier
in ranges that lead to
changes in the~measured signal yields 
as large as $\pm 20\%$. 
The~resulting difference 
in the~data\nobreakdash-simulation efficiency ratio 
is found to be~$3.0\%$. 

A~different class of systematic uncertainties 
directly affects the~fit itself, namely 
uncertainties associated 
%%
%% {\color{red}{\sout{to}}}
%% {\color{red}{
    with
%% }}
%%
the~fit 
models used to describe the~\jpsi\etapr\Kp, 
\psitwos\Kp and \jpsi\pip\pim spectra. 
The~systematic uncertainty is accounted for 
by fits with alternative models. 
The~list of alternative models to describe 
signal \Bp and \psitwos components includes 
a~modified Apollonios function~\cite{MartinezSantos:2013ltf},
which has exponential instead 
of power\nobreakdash-law tails, 
a~generalised 
Student's $t$-distribution~\cite{Student,Jackman} 
and 
a~modified Novosibirsk function~\cite{BaBar:2011qjw}.
For~the~combinatorial 
background for 
both the~signal and 
the~normalisation channels, 
a~third order polynomial function 
and an~exponential 
function multiplied by 
a~first order polynomial
function
are chosen 
as alternative models.
For~each fit only 
one component\,(either signal or background) 
is replaced at a~time, and 
the~same fit function is used for 
both signal and normalisation channel. 
The~largest deviation of the~ratio of signal
yields is found to be~0.6\%. 
As~alternative background models for \jpsi\pip\pim candidates, 
a~first order polynomial function and a~product of 
an~exponential function with a~first order polynomial 
function are used. 
The~largest deviation of 
the~signal yield is found to be 1.5\%. 
The~two deviations are added in quadrature 
to obtain a 1.6\% systematic uncertainty due to imperfect
knowledge of the signal and background shapes.


Finally, the~finite size of the~simulation samples 
contributes an~uncertainty of~$0.9\%$ 
on the~ratio of total efficiencies.
%%
The~total systematic uncertainty 
for the~ratio of branching fractions $\mathcal{R}$
is calculated as 
the~sum in quadrature of all the~values 
listed above and is found to be 6.0\%.

\begin{table}[tb]
\caption{Summary of systematic uncertainties
on 
the~ratio of branching 
fractions $\mathcal{R}$.
The~overall systematic uncertainty is calculated 
as a~sum 
in quadrature of all the sources.}
\label{tab:syst}
\begin{center}
\begin{tabular}{lc}
Source                           & Value [\%] 
\\[1mm]
\hline
\\[-2mm]
\Bp~kinematics 
& 0.1        \\
\Bu~decay model 
& 1.1        \\
Tracking efficiency correction   & 0.7        \\
Photon reconstruction correction & 3.6        \\
Kaon   identification            & 2.8        \\
Trigger efficiency               & 1.1        \\
Data-simulation agreement        & 3.0        \\
Fit model                        & 1.6        \\
Simulation sample size           & 0.9        
\\[1mm]
\hline
\\[-2mm]
Total & 6.0       
\end{tabular}\end{center} 
\end{table}


The~statistical significance 
for the~\mbox{$\decay{\Bu}{\jpsi\etapr\Kp}$}~decay
is recalculated using Wilks' theorem 
for each alternative fit model, 
and the~smallest values of 
$17$ and $10$~standard
deviations
for the~\etaprhog~and
\etapetapipi~cases, respectively, 
are taken as the~significance 
%% that accounts 
including 
the~systematic uncertainty.

\section{Results and summary}
\label{sec:sum}
The \betapk decays 
are observed for the first time
using 
proton\nobreakdash-proton collision data
collected by the~LHCb experiment  
at centre-of-mass energies of 
7, 8, and 13\tev,
corresponding 
to a~total integrated luminosity of $9\invfb$.
%% 
In~the~analysis, 
the~\etapr meson is reconstructed 
from the~two
$\rhoz\g$ and $\Peta\pip\pim$~final states.
The~signal significance 
exceeds 10 standard deviations
for both modes.
The~branching fraction of the~\betapk~decay 
is measured  for the~$\etaprhog$~sample 
through 
normalisation to 
the~known branching fraction of the~\bpsik~decay~\cite{PDG2022}.
The~ratio of branching fractions is found to be
\begin{equation*}
\dfrac{\BF(\betapk)}{\BF(\bpsik)} = 
\left(4.91\pm 0.47\pm0.29\pm0.07\right)\times10^{-2},
\end{equation*}
where the~first uncertainty is statistical, 
the~second  is systematic and the third  is related 
to 
the~uncertainties on 
the~branching fractions of 
%%
the~intermediate resonances.
%%
The~absolute branching fraction is determined using 
the~known branching fraction of \bpsik decays, 
\mbox{$\BF(\bpsik) = 
\left(6.24\pm0.20\right)\times10^{-4}$}~\cite{PDG2022},
and is found to be
\begin{equation*}
\BF(\betapk) = \left(3.06\pm0.29\pm0.18\pm0.04\right)\times10^{-5}\,,
\end{equation*}
where the~first uncertainty is statistical, 
the~second is systematic and 
the~third is due to external 
branching fractions uncertainties.
%%
%% {\color{red}{
    The~measured branching fraction
    is consistent with the~upper limit
    previously set by the~\belle collaboration~\cite{Belle:2006aam}.
%% }}
%%
An~inspection of the~$\jpsi\etapr$~mass spectrum 
shows no significant contributions from 
the~decays via intermediate   charmonium or 
charmonium\nobreakdash-like resonances. 
